(1) 小球内外静电势可设为均匀外场项与 \(r^{-2}\) 偶极项，匹配表面电势连续与法向位移矢量连续，得内部场 \(E_{\rm in}=3E_0/(\varepsilon_r+2)\)，偶极矩 \(p=4\pi\varepsilon_0a^3(\varepsilon_r-1)E_0/(\varepsilon_r+2)\)。因此 \(\boxed{\alpha_0=4\pi\varepsilon_0a^3(\varepsilon_r-1)/(\varepsilon_r+2)}>0\)。诱导偶极时间平均相互作用势为 \(U=-\operatorname{Re}\alpha\,|\boldsymbol E_0|^2/4=-\operatorname{Re}\alpha\,I/(2c\varepsilon_0)\)，梯度力 \(\boldsymbol F_{\rm grad}=-\nabla U=(\operatorname{Re}\alpha/(2c\varepsilon_0))\nabla I\)。系数 \(1/4\) 来自谐波时间平均与诱导偶极充电过程，不能直接使用永久静偶极能量。

(2) 对焦平面 \(I(r)=I_0e^{-2r^2/w^2}\)，有 \(\partial_r I=-4rI/w^2\)。因此 \(F_r=-[2\operatorname{Re}\alpha\,I_0/(c\varepsilon_0w^2)]r+O(r^3)\)，得 \(\boxed{\kappa_\perp=2\operatorname{Re}\alpha\,I_0/(c\varepsilon_0w^2)}\)，在 \(ka\ll1\) 下可取 \(\operatorname{Re}\alpha\simeq\alpha_0\)。若横向光势近似 \(U=\kappa_\perp x^2/2\)，Boltzmann分布或能量均分定理给 \(\boxed{\langle x^2\rangle=k_BT/\kappa_\perp}\)。这要求有热浴、阻尼与足够长的平衡时间；真空中无耗散粒子不自动热化。

(3) 令 \(\gamma=k^3\alpha_0/(6\pi\varepsilon_0)>0\)，则 \(\alpha=\alpha_0(1+i\gamma)/(1+\gamma^2)\)。所以 \(\operatorname{Im}\alpha=\alpha_0\gamma/(1+\gamma^2)=k^3|\alpha|^2/(6\pi\varepsilon_0)\)，满足无吸收散射的光学定理。对沿 \(+\hat z\) 传播的局部平面波，时间平均纵向力为 \(F_z=(k/2)\operatorname{Im}\alpha|\boldsymbol E_0|^2=k\operatorname{Im}\alpha\,I/(c\varepsilon_0)\)，故 \(\boxed{F_{\rm scatt}=k^4|\alpha|^2 I/(6\pi c\varepsilon_0^2)}\)。即使材料本身无吸收，辐射反作用仍使 \(\operatorname{Im}\alpha>0\)。

(4) 准静态 \(\alpha_0\propto a^3\)，故焦平面横向梯度力标度 \(F_{\rm grad}\sim a^3I/(cw)\)，而散射力 \(F_{\rm scatt}\sim k^4a^6I/c\)，二者比值约为 \(k^4a^3w\) 乘以介电对比因子。对于粒径增长或束腰很宽的光束，轴向散射推力可超过有限聚焦光束的轴向恢复梯度而使三维陷阱失稳；仅有横向稳定不等于完整三维束缚。当 \(ka\) 不再小、材料有吸收、场强引起非线性或布朗运动不平衡时，须使用Mie散射及非平衡随机动力学。
